\section{Week 8}

Recall that a lot of the problems given in the homework are based on more general theory of the results that we have been going thus far. For instance, one can prove that 
If \( G  \) is cylic and \( \varphi  \) is a homomorphism from \( G  \) to \( H  \), then \( \varphi(G)  \) is cyclic. Note that \( \pi : G \to G / N  \) where \( g \mapsto gN \)   and \( \pi  \) is a homomorphism. Another problem we can do is prove that if \( G  \) is cyclic, then \( G / N  \) is cyclic. 

For an \( \psi  : G \to H  \) where \( \psi  \) is a group homomorphism. If \( | x  |  =  n < \infty   \), then \( | \psi(x)  | | x  |  \). Because \( \pi  \) is a homomorphism from \( G \to G / N  \). Hence, the homework problem 2 follows immediately.

Last time we saw that for \( H , K \leq G  \) we have 
\begin{align*}
    | H K  |  &= \frac{ | K  |  | H  |  }{  | K \cap  H |  }.
\end{align*}
Note that \( HK  \) may always be a subgroup of \( G  \).

\begin{eg}
   Let \( {S}_{3} = G  \) and \( H = \langle (1 \ 2) \rangle \) and \( K = \langle  (2 \ 3) \rangle \). Then 
   \[  | H K |  = \frac{ | H  |  | K  |   }{ | H \cap K  |  }  = \frac{ 2 \cdot 2  }{ 1  }  = 4.  \]
   However, we see from Lagrange's theorem that if \( H K \leq G  \), then \( 4 \mid 3!  \) which is a contradiction. We can further deduce that 
   \[  \langle (1  \ 2 )  , (1 \  3) \rangle = {S}_{3}. \]
   Since \( 4 \leq | \langle (1  \  2 ) , (2 \ 3) \rangle | \leq 6  \) and must also divide \( 6  \). So, \( (1 \ 2) \) and \( (2 \ 3 ) \) generates all of \( {S}_{3} \).
\end{eg}

\begin{prop}
   If \( H, K  \) are subgroups of a group \( G  \), \( H K \leq G  \) if and only if \( HK =KH  \).
   \textit{Note that \( HK = KH  \) does not indicate that the elements of the set commute with each other. Only that \( h k  \in H K \), we have that \( h k = {k }_{1} {h}_{1}  \) for some \( {h}_{1} \in H  \) and \( {k}_{1} \in K  \)} 
\end{prop}

\begin{proof}
Assume that \( HK = KH  \). Our goal is to show that \( HK  \) is a subgroup of \( G  \). Since \( H  \) and \( K  \) are non-empty, then \( HK \neq \emptyset \). It remains to show that given \( a,b \in H K  \) that \( a b^{-1} \in HK \). To this end, let \( a,b  \in H  K  \). Then \( a = {h}_{1} {k}_{1} \) and \( b = {h}_{2} {k }_{2} \). Then 
\begin{align*}
    a b^{-1} &= ({h}_{1} {k}_{1}) ({h}_{2} {k}_{2})^{-1} \\
             &= ({h}_{1} {k}_{1}) ({k}_{2}^{-1} {h}_{2}^{-1}) 
\end{align*}
Since \( K \leq G  \), we have \( {k}_{1} {k}_{2}^{-1} = {k}_{3} \in K  \) and \( {h}_{2}^{-1} = {h}_{3} \in H \). This gives that 
\[ a b^{-1} = {h}_{1} {k}_{3} {h}_{3}. \]
Since \( H K = K H \), we know that \( {k}_{3} {h}_{3} \in HK \). That is, \( {k}_{3} {h}_{3} = {h}_{4} {k}_{4}  \) for some \( {h}_{4} \in H  \) and \( {k}_{4} \in K  \). So, 
\[  a b^{-1} = {h}_{1} {h}_{4} {k}_{4} \]
and letting \( {h}_{1} {h}_{4} = {h}_{5} \) which is in \( H  \) by closure, we have that 
\[  a b^{-1} = {h}_{5} {k}_{4} \in HK.  \]

Conversely, suppose that \( H K \leq G  \). Our goal is to show that \( HK = K H  \); that is, \( HK \subseteq  KH   \) and \( KH \subseteq  HK \). Since \( K \leq HK  \) and \( H \leq HK  \) and so \( KH \leq HK \) by closure of \( H K  \). To show that \( HK \leq KH  \), let \( hk \in HK \). Since \( HK  \) is a subgroup, then \( hk  \) is the inverse of some element \( a \in HK \). That is, \( h k = a^{-1} = ({h}_{1} {k}_{1})^{-1} = {k}_{1}^{-1} {h}_{1}^{-1} \in KH \). It follows that \( HK  \leq KH \). Furthermore, \( HK = KH \). 
\end{proof}


\begin{eg}
    Let \( G = {D}_{8}  \), \( H = \langle  r  \rangle \) and \( K  = \langle   s  \rangle \). Notice that \( rs \in H K \) and \( rs = s r^{-1} \in KH  \). By the way, 
    \[ | H K |  = \frac{ | H  |  | K  |  }{ | H \cap  K |   }  = \frac{ 4 \cdot 2  }{  1  }  = 8.  \] 
    So, \( HK = {D}_{8}  = KH  \).
\end{eg}

\begin{corollary}
If \( H  \) and \( K  \) are subgroups of a group \( G  \) and \( H \leq {N}_{G}(K) \), then \( H K  \) is a subgroup of \( G  \). In particular, if \( K \trianglelefteq G  \), then \( H K  \) is a subgroup of \( G  \) for all \( H \leq G  \). \textit{Note: If \( K \trianglelefteq G  \), then \( {N}_{G}(K) = G  \) and certainly \( H \subseteq  {N}_{G}(K) \)   for any \( H \leq G  \)}
\end{corollary}

\begin{proof}
We will prove that \( HK = KH  \). Let \( h \in H  \) and \( k \in K  \). By assumption, we have that \( h \in {N}_{G}(K) \); that is, \( h k h^{-1} \in K  \). Hence, \( h k = h k (h^{-1} h) \) where \( h k h^{-1} \in K  \) and \( h \in H  \). Thus, \( hk \in KH  \). Thus, \( HK \subseteq  KH  \). Similarly, we want to show that \( kh  = h h^{-1} kh  \) where \( h \in H  \) and \( h^{-1} k h  \in  K \) since \( h \in {N}_{G}(K) \) implies that \( h^{-1} \in {N}_{G}(K) \). It follows that \( KH \subseteq  HK  \) and hence \( HK = KH  \).
\end{proof}

\begin{theorem}[ ]
    Let \( H,K  \) be subgroups of a group \( G  \) with \( H \leq K \leq G  \). Then the index 
    \[  | G : H  |  = | G : K  | | K : H  |. \]
\end{theorem}
\begin{proof}
Let \( {g}_{i} \) be a distinct representative for a left coset of \( H  \) in \( G  \) for all \( i \in I  \) where \( I  \) is some indexing set. So, we have  
\[  \{  {g}_{i} H  :  i \in I  \}  = G / H = \{  g H : g \in G  \}.   \]
We have \( {g}_{i} H = {g}_{j} H  \) if and only if \( {g}_{i} = {g}_{j} \) (we are counting the elements of \( g H  \) because we don't know if it will be finite or infinte). \textit{At this point we are invoking the axiom of choice.} Let \( \psi : I \times K / H \to G / H  \) by 
\[  \psi(i, kH) = {g}_{i} kH.  \]
To see that this is well-defined, let us suppose that \( {k}_{1} H = {k}_{2} H  \) for some \( {k}_{1}  \) and \( {k}_{2}  \) in \( K  \). That is, \( {k }_{1}^{-1} {k}_{2} \in H \). Then \( \psi (i, {k}_{1}H) = {g}_{i} {k}_{1} H  \) and \( \psi(i, {k}_{2} H ) = {g}_{i} {k}_{2} H  \). So, 
\begin{align*}
    ({g}_{i} {k}_{1})^{-1} ({g}_{i} {k}_{2}) &= ({k}_{1}^{-1} {g}_{i}^{-1})({g}_{i} {K}_{2}) \\
                                             &= {k}_{1}^{-1} {k}_{2} \in H 
\end{align*}
by assumption. Hence, \( \psi  \) is well-defined. We now show that \( \psi  \) defines a bijection. Suppose that 
\begin{align*}
&\psi(i, {k}_{1} H ) = \psi (j, {k}_{2} H ) \\
&\implies ({g}_{i}{k}_{1}) H = ({g}_{j} {k}_{2}) H \\
&\implies ({g}_{i} {k}_{1})^{-1} ({g}_{j} {k}_{2}) \in H. 
\end{align*}
This implies that \( {k}_{1}^{-1} {g}_{i}^{-1} {g}_{j} {k}_{2} = h   \) (call this (*)) for some \( h \in  H \). This tells us that \( {g}_{i}^{-1} g_{j} = {k}_{1} h {k}_{2}^{-1} \) for some \( h \in H  \). Note that \( H \leq K \leq G  \) and so \( {g}_{i}^{-1} {g}_{j} \in K  \) since \( H \subseteq  K  \). Then \( {g}_{i} K = {g}_{j} K  \). Thus, we have \( {g}_{i} = {g}_{j} \). This further implies that \( i = j  \). Using this in (*) gives 
\begin{align*}
    &{k}_{1}^{-1} {h}_{i}^{-1} {g}_{j} {k}_{2} = h  \\
    &\implies {k}_{1}^{-1} {k}_{2} = h \in H \\
    &\implies {k}_{1} H = {k}_{2} H.
\end{align*}
Hence, we see that \( \varphi  \) is one-to-one. Let \( g H \in G / H  \). Since the left cosets of \( k  \) partition the group \( G  \), we have that \( g \in {g}_{i} K  \). That is, \( g = {g}_{i} k  \) for some \( k \in K  \). Hence, \( \varphi(i, k H ) = {g}_{i} k H = g H  \). Thus, \( \psi  \) is surjective. Hence, 
\[  \underbrace{| G : H  |}_{G /H }  = \underbrace{| G : K  |}_{I}  \cdot \underbrace{| K : H  |}_{K/H}. \]
\end{proof}

\begin{theorem}[First Isomorphism Theorem]
    If \( \varphi : G \to H  \) is a homomorphism of groups, then 
    \[  \text{ker}(\varphi) \trianglelefteq G, \]
    and
    \[  G / \text{ker}(\varphi) \cong \varphi(G). \]
\end{theorem}
\begin{proof}
\textit{Proof is left to the reader.} Take \( \mu: G / \text{ker}(\varphi) \to \varphi(G) \) where \( g K \mapsto \varphi(g) \) and prove that is well-defined and is a bijection. 
\end{proof}

